\documentclass[11pt]{article}
\usepackage{amsmath,amssymb}
\title{Interior background shifts: valid relative-box identity and its limit}\author{}\date{}
\begin{document}\maketitle
\paragraph{Primary passage checked.}
L\"uscher--Weisz, \emph{Background field technique and renormalization
in lattice gauge theory}, hep-lat/9504006, section6.2,
equations6.13--6.24, is available at
\texttt{https://arxiv.org/pdf/hep-lat/9504006}.
Their lattice split shift preserves the total link and commutes with
BRS on ghost-independent functions. The gauge-fixing/ghost variation
is BRS exact. Their source-free $W[B,0]$ assertion is not an assertion
that a Legendre effective action or fixed-boundary free energy is zero.
No infinite-volume renormalizability statement is imported below.

\paragraph{1. Exact interior-supported relative-box shift.}
Write the actual total link as $U_b=Q_bV_b$, with $V$ an auxiliary
background and $Q=\exp(\alpha z)$. For a deterministic Lie-algebra
field $\omega$ supported ONLY on integrated interior links, set
\[
 V_b(t)=e^{\alpha t\omega_b}V_b,\qquad
 Q_b(t)=Q_b e^{-\alpha t\omega_b}. \tag{1}
\]
Then $U_b(t)=U_b$ exactly. All original boundary links and their
constraints stay fixed. Right translation preserves the original
link Haar measure and has zero Haar divergence.
Thus for any actual observable $O(U)$ and physical cutoff $\chi(U)$,
\[
 \int dQ\,\chi(QV)e^{-S(QV)}O(QV)
 \quad\hbox{is unchanged by the interior split shift.} \tag{2}
\]
This is an exact finite group-variable identity with the original
raw scalar. It changes the split, not the physical total field.

For the actual transported internal source of the internal source tree to fp transport calculation and the two stage quotient and compensated covariance calculation,
\[
 O_t(U)=\exp\{t(Aw)^*q_{\rm tree}(U)-t^2w^*Aw/2\}, \tag{3}
\]
the shift variation is ZERO because the full original tree quotient
depends only on $U$. Its physical-chart cutoff is likewise unchanged.
Replacing $q_{\rm tree}$ by the bare quantum logarithm would invalidate
this statement. The nonlinear source contacts of the internal source tree to fp transport calculation are required.

\paragraph{2. Ghost boundary and the source-current Ward.}
In the relative local FP formulation, ghosts are assigned only to
integrated vertices and are zero at the original fixed boundary.
An interior shift changes no such domain. On a ghost-independent
observable, BRS acts through the total link
$sU_b=c_sU_b-U_bc_t$, whereas the split shift acts at fixed $U$.
The two variations therefore commute directly, without integration
by parts across a physical boundary.

Write the local gauge-fixing and ghost part as $s\Psi$ (or use the
equivalent on-shell expression of the cited passage), and set
$R=\delta_s\Psi$, which is odd. Then
$\delta_sS_{\rm tot}=sR$. With source $J$ and cutoff $\chi$,
the formal local BRS integration identity gives
\[
 \delta_s\log Z[J]=
 \left\langle\delta_sJ-RsJ
       +\chi^{-1}\delta_s\chi+\chi^{-1}(s\chi)R\right\rangle_J
       +{\cal B}_{\rm domain}. \tag{4}
\]
The domain term is the actual local FP-chart integration flux; it
cannot be erased by asserting a global nonperturbative BRS measure.
Equation (4) follows by applying $s$ to $\chi R e^J$ with its odd
product sign. Fixed-volume formal Laplace coefficients may use the
original fixed physical chart $\chi_{\rm phys}(\alpha q)$ and the
same local gauge-noise construction as the local gauge completion first nonlinear coefficient calculation. No cutoff bound uniform
in volume is proved here.

For (3) and its gauge-invariant physical cutoff, $sJ=\delta_sJ=0$
and $s\chi=\delta_s\chi=0$. The surviving content of (4) is exactly
the representation independence in (2), with any local-chart flux
handled on its original domain. It is not a new physical coercivity
or mean-current estimate.

\paragraph{3. No equation-of-motion gain for the completed covariance.}
The source covariance in the two stage quotient and compensated covariance calculation is a derivative in the physical observable
parameter $t$, not the auxiliary split parameter of (1).
The split identity is valid for every $t$, so it transports all its
derivatives consistently; it does not force any of them to vanish.
In particular it does not imply
$\operatorname{Var}((Aw)^*q_{
m tree})=w^*Aw$.
The actual $R=2$ values recorded in the two stage quotient and compensated covariance calculation,
$K_{\rm tree}=-43/24$ or $-47/24$, already exclude that inference.

The remaining completed tensor is still
\[
 -[\alpha^2]\left\{E_{\nu_\alpha}v_\alpha
       +\operatorname{Var}_{\nu_\alpha}m_\alpha-w^*Aw\right\}, \tag{5}
\]
with the original conditional normalizers, transported tree source,
and induced coarse exchange of the same coefficient schur cumulant decomposition calculation. Its harmonic part is controlled;
the residual/coarse response and heat-regulator source contacts are
not. The background shift supplies no extra long-distance or energy
factor in (5). A genuine new estimate of this completed covariance
or its equivalent signed current kernel is still required.
This particular shift route is therefore exhausted as a quantitative
advance: it validates transport but does not close the main norm.
The original nonlinear remainder and S94 remain open.
\end{document}
